\chapter{Hedging and Inventory in Practice}\label{ch:hf:hedging-and-inventory-in-practice}

At ten to four a stock market maker is long a basket of forty names it never chose. Its clients sold into a falling afternoon, its quotes leaned
against each position as the models of chapters 3 and 4 say they should, and still the book holds \$1.6 million of stock, \$830,000 of it in
index-equivalent exposure. It can sell the stocks, sell the index future against them, or keep them overnight and start tomorrow long. In this
chapter's simulated market, selling the stocks costs \$440; selling the future costs \$83 and leaves an overnight standard deviation of \$3,100;
keeping the book leaves a one-in-a-hundred overnight loss of \$14,800. During the day the same question arises every few seconds, and answering it
at every move costs 14\% of the spread the desk earns.

\section{Choosing the hedge instrument}

A market maker's inventory is a by-product: it holds what clients sold to it. The inventory models of chapters 3 and 4 manage it through the quotes,
which cost nothing but act slowly. A hedge acts at once but crosses a spread.

\begin{definition}[Hedge instrument]\label{def:hf:hedging-and-inventory-in-practice:instrument}
A \emph{hedge instrument}\index{hedge instrument} is a security a market maker trades, not to earn a spread but to offset the risk of positions it
acquired by making markets; it is chosen for the cost of trading it and for how much of the positions' risk it removes.
\end{definition}

The choice is between instruments that remove different parts of the risk at different prices. In a one-factor model, name $i$ has price $S_i$, a
beta $\beta_i$ to the market factor and an idiosyncratic daily volatility $\sigma_i$; the factor has daily volatility $\sigma_F$ and the index future
tracks it.

\begin{definition}[Delta-equivalent inventory]\label{def:hf:hedging-and-inventory-in-practice:delta}
The \emph{delta-equivalent inventory}\index{delta-equivalent inventory} of a book is its exposure expressed in one reference instrument: each
position's value times its sensitivity to the reference, summed across instruments. Against an index, a book of stocks $q_i$ has
\[
D=\sum_i q_i\,S_i\,\beta_i
\]
dollars of index exposure; an ETF counts with its own beta, an option with its delta times its underlying's beta.
\end{definition}

The book's daily variance is $D^2\sigma_F^2+\sum_i (q_iS_i\sigma_i)^2$. Selling $D$ dollars of the future removes the first term at a cost $c_F|D|$,
with $c_F$ the future's cost per dollar traded; trading every stock back to zero removes both at a cost $c_S\sum_i|q_iS_i|$. With a charge $\lambda$
on each squared dollar of variance held for a horizon $h$, the choice compares three numbers (\cref{lst:hf:hedging-and-inventory-in-practice:choose}):
\[
\lambda h\,\bigl(D^2\sigma_F^2+\textstyle\sum_i(q_iS_i\sigma_i)^2\bigr),\qquad c_F|D|+\lambda h\textstyle\sum_i(q_iS_i\sigma_i)^2,\qquad c_S\textstyle\sum_i|q_iS_i|.
\]
The future wins when the book is concentrated in the factor and large, the stocks when the idiosyncratic risk is large or the horizon long, and doing
nothing when the book is small. A future also brings risks the one-factor model omits: its basis to the cash index, the book's betas estimated with
error, and, for a book of stocks outside the index, a factor exposure the future does not match. Chapter 9 of One Quant Book 9 treats the
estimated hedge ratio.

\section{Partial hedging and hedge triggers}

Hedging is rarely all or nothing. If the only cost is proportional, the one-period problem has a closed form.

\begin{proposition}[The dead zone]\label{prop:hf:hedging-and-inventory-in-practice:dead}
Selling $x$ dollars of the future against an exposure $D$ at cost $c_F|x|$ with a risk charge $\lambda\sigma_F^2(D-x)^2$ is optimal at
\[
x^\ast=\operatorname{sign}(D)\,\max\Bigl(|D|-\frac{c_F}{2\lambda\sigma_F^2},\,0\Bigr):
\]
the market maker hedges down to the edge of a dead zone of half-width $c_F/(2\lambda\sigma_F^2)$ and not inside it.
\end{proposition}

\begin{proof}
For $0<x<D$ the derivative of $c_Fx+\lambda\sigma_F^2(D-x)^2$ is $c_F-2\lambda\sigma_F^2(D-x)$, zero at $D-x=c_F/(2\lambda\sigma_F^2)$; the objective is
convex, and no hedge beyond $D$ or of the opposite sign helps. If $|D|$ is inside the half-width, the derivative at $x=0$ is already non-negative.
\end{proof}

Over a day, the exposure moves continuously as fills arrive, and the one-period answer becomes a band: do nothing while the unhedged exposure is
inside $\pm b$, trade back to the edge when it leaves. One Quant Book 4, chapter 10 derives the optimal band for an exposure that diffuses with
volatility $\sigma_D$, a running risk charge $\gamma X^2$ and a proportional cost $c$: $b^\ast=\bigl(3c\sigma_D^2/(4\gamma)\bigr)^{1/3}$. Here
$\gamma=\lambda\sigma_F^2$. The cube root makes the band robust: doubling the cost widens it by 26\%.

A hedge trigger may also be an event rather than a level: a large fill, a fill from a client known to be informed, or a move in the market that
makes the exposure's risk jump. The discipline is the same: each hedge must be worth more in risk removed than it costs.

\section{Inventory across correlated books}

A firm that makes markets in the index's stocks, its ETFs and its futures holds one factor exposure spread across books. Netting the
delta-equivalent inventories before hedging is the cheapest hedge there is: a long in one book offsets a short in another, and the firm hedges only
the sum. The desk that makes markets in ETFs (chapter 12) can often hedge with its own stock inventories; the options desk's delta (chapter 19) adds
to the same number.

Netting has two conditions. The books must be marked and summed in near real time, which is an engineering problem; and the residual risks that
netting leaves (idiosyncratic, basis, time zone) must stay inside each book's own limits.

The empirical record says inventory matters to prices and to liquidity. In twelve years of daily New York Stock Exchange specialist data from 1994 to
2005, Hendershott and Menkveld measured price pressures from intermediaries' inventory of 49 basis points on average, with a half life of 0.92
days: 17 basis points and 0.54 days in the largest stocks, 118 basis points and 2.11 days in the smallest. Comerton-Forde, Hendershott, Jones,
Moulton and Seasholes found in eleven years of the same specialists' data that market-wide and firm-level spreads widen when specialists hold large
positions or lose money, most strongly when inventories are big or results poor. And the high-frequency market maker Menkveld studied in 2013 earned
a profit on the spread and lost on its positions: of \euro1.55 a trade earned on the spread net of fees, \euro0.68 went on positioning.

\section{End of day and overnight}

\begin{definition}[End-of-day flattening]\label{def:hf:hedging-and-inventory-in-practice:flatten}
\emph{End-of-day flattening}\index{end-of-day flattening} is the reduction of a market maker's positions towards zero before the close, through its
quotes, through orders in the continuous market or in the closing auction, or by hedging what remains with an instrument that trades overnight.
\end{definition}

The night changes the problem: the stocks cannot be traded until the open, the gap at the open is not a diffusion but a jump with fat tails, and
the book has no spread income to set against it. Four ways to flatten, from cheapest to dearest:

\begin{itemize}
\item \emph{Through the quotes.} Skew harder in the last hour. In the tutorial, quotes skewed four times harder from 15:00 halve the gross book at
the close, \$1.32 million to \$670,000, while the desk trades the same number of fills.
\item \emph{In the closing auction.} No spread to cross, a single price; the risk is the auction's own imbalance, and a large order moves the close
(One Quant Book 1, chapter 13).
\item \emph{In the continuous market} before the close: the half-spread plus impact.
\item \emph{Keep the stocks and sell the future overnight}: the factor risk goes, the idiosyncratic gaps stay.
\end{itemize}

\begin{omfigure}
\begin{tikzpicture}
\begin{axis}[width=11cm,height=6.4cm,xlabel={\small intraday risk (\$ thousand a day)},ylabel={\small hedging cost (\$ thousand a day)},
  xmin=0,xmax=10,ymin=-1,ymax=14.5,tick label style={font=\footnotesize},grid=major,grid style={gray!25},table/col sep=comma]
\addplot[omThm,only marks,mark=*,mark size=1.8pt] table[x=risk,y=cost]{figdata/hft/11-hedging-and-inventory-in-practice/frontier.csv};
\begin{scope}[font=\scriptsize]
\node[anchor=west] at (axis cs:2.7,12.76) {stocks};
\node[anchor=west] at (axis cs:4.6,6.70) {future at every move};
\node[anchor=east] at (axis cs:4.35,3.14) {band \$100k};
\node[anchor=west] at (axis cs:4.9,2.3) {band \$250k};
\node[anchor=north] at (axis cs:5.44,0.6) {band \$500k};
\node[anchor=south] at (axis cs:6.99,0.6) {band \$1m};
\node[anchor=south] at (axis cs:8.9,0.3) {none};
\end{scope}
\end{axis}
\end{tikzpicture}

\omcaption{Hedging cost against intraday risk (the one-minute P\&L's standard deviation scaled to a day) for seven rules: no hedge, the future kept
within half a contract of the book's \omterm{def:hf:hedging-and-inventory-in-practice:delta}{delta-equivalent inventory}, the future outside no-trade bands of \$100,000 to \$1 million, and every name traded
back to \$20,000 in its own stock. Forty names, ten simulated days, spread income \$46,800 a day. Data:
\texttt{hf\_hedging.frontier}.}\label{fig:hf:hedging-and-inventory-in-practice:frontier}
\end{omfigure}

\section{Tutorial: forty names, one factor, one future}\label{tut:hf:hedging-and-inventory-in-practice:tut}

\begin{tutorial}
\textbf{Goal.} Measure what each hedging rule costs and what risk it leaves on a simulated day of fills across forty names, then take the ten-to-four
decision. \textbf{End state:} \cref{fig:hf:hedging-and-inventory-in-practice:frontier,fig:hf:hedging-and-inventory-in-practice:overnight} and the table
of the ten-to-four decision.
\begin{tutsteps}
\item \textbf{The day.} \texttt{firm.mmhedge.FlowDay} draws forty names (prices \$20--150, betas 0.6--1.4, idiosyncratic volatility 1--2\% a day,
factor 1\%), 0.05 fills a second per name of 100 shares at a half-spread of one cent, quotes skewed against inventory, and clients who sell after the
factor falls (a tilt of 0.15 per standard deviation of the last minute's move). The tilt makes the names' inventories move together: the book's
delta-equivalent exposure has a daily volatility of \$2.5 million. Simulating fills rather than forty full order books keeps the day at a fraction of
a second; chapter 1's harness is the tool when fills depend on the book.
\item \textbf{The choice} of instrument (\cref{lst:hf:hedging-and-inventory-in-practice:choose}).
\omcode{code/firm/mmhedge/firm_mmhedge.py}{47}{56}{Cost plus the charge on the variance left, for the three choices.}\label{lst:hf:hedging-and-inventory-in-practice:choose}
\item \textbf{The rules} on the future: within half a contract (\$50,000 notional), or a no-trade band.
\omcode{code/firm/mmhedge/firm_mmhedge.py}{125}{144}{The future position along the day's exposure, and what it cost.}\label{lst:hf:hedging-and-inventory-in-practice:band}
\item \textbf{The frontier.} Ten days, seven rules, future at half a basis point a dollar, stocks at 1.7 basis points (\cref{fig:hf:hedging-and-inventory-in-practice:frontier}).
\item \textbf{Ten to four.} Take the book at 15:50 of the first day and compare flattening in the last ten minutes (half-spread plus square-root
impact), selling the future overnight (in and out: one basis point) and keeping it, with Student-$t$ gaps of four degrees of freedom (factor 0.6\%,
names 1\%).
\end{tutsteps}
\textbf{What to change next.} Add a basis between the future and the factor; estimate the betas from a month of simulated returns instead of using
the true ones; net a second book (an ETF) into $D$.
\end{tutorial}

The frontier (\cref{fig:hf:hedging-and-inventory-in-practice:frontier}): with no hedge, intraday risk is \$8,540 a day. Keeping the future within
half a contract halves it to \$4,410 but trades 2,682 contracts and costs \$6,700 a day, 14\% of the spread income. A band of \$250,000 leaves
\$4,760 of risk for \$1,880. The stocks remove the idiosyncratic risk too (\$2,525) for \$12,760. With a charge of $\lambda=10^{-4}$ per squared
dollar, the band of \$250,000 scores best (\$42,570 of spread income less costs and charge, against \$39,640 unhedged and \$38,260 for the future at
every move), and the formula's band is \$284,000.

\begin{omfigure}
\begin{tikzpicture}
\begin{axis}[width=11cm,height=5cm,xlabel={\small overnight loss (\$ thousand)},ylabel={\small probability},xmin=-25,xmax=35,ymin=0,
  yticklabel style={/pgf/number format/fixed,/pgf/number format/precision=2},scaled y ticks=false,
  tick label style={font=\footnotesize},grid=major,grid style={gray!25},table/col sep=comma,
  legend style={font=\scriptsize,at={(0.98,0.95)},anchor=north east}]
\addplot[omThm,thick,const plot mark mid] table[x=loss,y=keep]{figdata/hft/11-hedging-and-inventory-in-practice/overnight.csv};
\addplot[omDef,thick,dashed,const plot mark mid] table[x=loss,y=future]{figdata/hft/11-hedging-and-inventory-in-practice/overnight.csv};
\draw[omThm,dotted] (axis cs:14.78,0) -- (axis cs:14.78,0.06);
\draw[omDef,dotted] (axis cs:7.36,0) -- (axis cs:7.36,0.06);
\legend{kept,{kept, future sold}}
\end{axis}
\end{tikzpicture}

\omcaption{Overnight loss of the book held at 15:50 (\$1.57 million gross, \$834,000 delta-equivalent), kept as it is or with \$834,000 of future
sold, in bins of \$1,000; Student-$t$ gaps with four degrees of freedom, factor 0.6\% and names 1\% a night; 200,000 draws. Dotted lines: the 99\%
quantiles, \$14,780 and \$7,360. Data: \texttt{hf\_hedging.ten\_to\_four}.}\label{fig:hf:hedging-and-inventory-in-practice:overnight}
\end{omfigure}

\begin{center}
\small
\begin{tabular}{@{}lrrrr@{}}
\toprule
at 15:50 & cost & overnight s.d. & 99\% loss & 99\% expected shortfall\\
\midrule
flatten in the last ten minutes & \$438 & 0 & 0 & 0\\
sell the future overnight & \$83 & \$3,086 & \$7,358 & \$8,813\\
keep & 0 & \$5,879 & \$14,776 & \$19,795\\
\bottomrule
\end{tabular}
\end{center}

With the same charge, the three choices score \$438, \$1,035 and \$3,454: flattening wins, because the idiosyncratic gaps the future cannot remove
are charged at a night's variance with no spread to earn against them. Skewing the quotes from 15:00 would have halved the flattening bill (\$163
against \$381 on average over the ten days).

\section{Build: the hedging module}\label{bld:hf:hedging-and-inventory-in-practice:mmhedge}

\begin{build}
\bfield{Purpose} Express a market maker's positions as one factor exposure, choose what to hedge with, and apply band, trigger and end-of-day rules.
\bfield{Interface} \texttt{delta\_equivalent}, \texttt{book\_var}, \texttt{choose(q, p, beta, sf, se, c\_fut, c\_stk, lam, horizon)},
\texttt{partial\_hedge}, \texttt{optimal\_band} (from \texttt{firm.impulse}), \texttt{FlowDay(n, seed)} with
\texttt{inventory(stock\_band, c\_stk, eod\_from, eod\_skew)}, \texttt{hedge\_future(D, rule, band, notional, c)}, \texttt{day\_pnl},
\texttt{overnight}, \texttt{flatten\_cost}.
\bfield{Rules} Costs per dollar traded; the future trades in whole contracts; the band's edge, not zero, is the target after a trigger; overnight
gaps are fat-tailed.
\bfield{Acceptance tests} \texttt{code/firm/mmhedge/tests/}: exposure and variance by hand and the choice in three constructed cases; the dead zone is
the numerical argmin; the future rules keep the exposure within their bounds and the band trades a fifth as much; the tilt makes names co-move and the
stock band caps every name; overnight standard deviation matches the formula; flattening cost by hand; the band matches the closed form; a harder
skew from 15:00 halves the book at the close.
\bfield{Stretch} Basis and beta estimation error; several books netted; a closing-auction flattening with imbalance risk.
\end{build}

\begin{omsources}
\item T.~Hendershott, A.~J.~Menkveld, Price pressures, \emph{Journal of Financial Economics} 114(3), 2014, 405--423.
\item C.~Comerton-Forde, T.~Hendershott, C.~M.~Jones, P.~C.~Moulton, M.~S.~Seasholes, Time variation in liquidity: the role of market-maker
inventories and revenues, \emph{Journal of Finance} 65(1), 2010, 295--331.
\item A.~J.~Menkveld, High frequency trading and the new market makers, \emph{Journal of Financial Markets} 16(4), 2013, 712--740.
\end{omsources}

\section{Exercises}

\begin{exercise}[$\star$]\label{exo:hf:hedging-and-inventory-in-practice:1}
A book is long 1,000 shares of a \$50 stock with beta 1.2 and short 500 shares of a \$100 stock with beta 0.8. What is its \omterm{def:hf:hedging-and-inventory-in-practice:delta}{delta-equivalent inventory},
and how many \$50,000 futures does the half-contract rule sell?
\end{exercise}

\begin{exercise}[$\star$]\label{exo:hf:hedging-and-inventory-in-practice:2}
With $c_F=0.5$ basis point, $\lambda=10^{-4}$ and $\sigma_F=1\%$, what is the dead zone's half-width, and how much of an \$833,747 exposure is hedged?
\end{exercise}

\begin{exercise}[$\star$]\label{exo:hf:hedging-and-inventory-in-practice:3}
What does it cost to sell \$833,747 of future at 15:50 and buy it back at the open, at 0.5 basis point each way?
\end{exercise}

\begin{exercise}[$\star\star$]\label{exo:hf:hedging-and-inventory-in-practice:4}
The future removes the factor risk. Why does the hedged book still have an overnight standard deviation of \$3,086?
\end{exercise}

\begin{exercise}[$\star\star$]\label{exo:hf:hedging-and-inventory-in-practice:5}
Compute the optimal band at $\lambda=10^{-5}$, $3\times10^{-5}$ and $10^{-4}$ with $\sigma_D=\$2.476$ million a day, and explain why the band
narrows as risk aversion rises.
\end{exercise}

\begin{exercise}[$\star\star$]\label{exo:hf:hedging-and-inventory-in-practice:6}
Skewing the quotes harder from 15:00 halves the book at the close. Why is it nearly free, and when would it not be?
\end{exercise}

\begin{exercise}[$\star\star\star$]\label{exo:hf:hedging-and-inventory-in-practice:7}
\emph{Coding.} Double the future's cost (\texttt{frontier(1e-4)}). Which simulated band scores best, and what does the formula say?
\end{exercise}

\begin{exercise}[$\star\star\star$]\label{exo:hf:hedging-and-inventory-in-practice:8}
\emph{Find the flaw.} ``We sell a future every time the book's beta moves by one contract. That minimises our risk, so it must be the right policy.''
\end{exercise}

\section{Problem: Forty Names at Ten to Four}

\begin{problem}[{Weekend problem --- forty names at ten to four}]\label{pb:hf:hedging-and-inventory-in-practice:1}
A stock market maker ends its day long a basket it did not choose and must decide how to hedge during the day and what to hold overnight.

\medskip\noindent\textbf{Part I --- The exposure.}
\begin{enumerate}
\item Define a \omterm{def:hf:hedging-and-inventory-in-practice:instrument}{hedge instrument} and the costs that choose it.
\item Define the \omterm{def:hf:hedging-and-inventory-in-practice:delta}{delta-equivalent inventory} of a book of stocks against an index.
\item Write the book's daily variance in the one-factor model.
\item Why do clients who sell into a falling market make a market maker's names co-move?
\end{enumerate}

\medskip\noindent\textbf{Part II --- Rules.}
\begin{enumerate}[resume]
\item State and prove the dead zone (\cref{prop:hf:hedging-and-inventory-in-practice:dead}).
\item Give the optimal band and its value at $\lambda=10^{-4}$.
\item Name three risks of a future hedge that the one-factor model omits.
\item Why net across books before hedging, and on what conditions?
\end{enumerate}

\medskip\noindent\textbf{Part III --- Measurements.}
\begin{enumerate}[resume]
\item Give the cost and intraday risk of the four main rules.
\item Which rule scores best at $\lambda=10^{-4}$, and how close is the formula?
\item What did the empirical studies find about inventory and prices, and inventory and spreads?
\item What did Menkveld's market maker earn on the spread and lose on positions?
\end{enumerate}

\medskip\noindent\textbf{Part IV --- The verdict.}
\begin{enumerate}[resume]
\item State the \emph{named result}: the cost and residual risk of the three choices at 15:50, and the overnight loss distribution of keeping the book.
\item Which choice wins with the same charge, and why?
\item Define \omterm{def:hf:hedging-and-inventory-in-practice:flatten}{end-of-day flattening} and list its four routes.
\item What does skewing from 15:00 change in the flattening bill?
\item When would keeping the book overnight be right?
\item How does the answer change if the betas are estimated with error?
\item What would an options book add to the \omterm{def:hf:hedging-and-inventory-in-practice:delta}{delta-equivalent inventory}?
\item In one sentence: when is a hedge worth it?
\end{enumerate}
\end{problem}

\section{Interview questions}

\begin{interviewq}[$\star$ \iqroles{trader}]\label{iq:hf:hedging-and-inventory-in-practice:1}
You are long \$5 million of bank stocks at 15:55. What are your options, and which do you take?
\end{interviewq}

\begin{interviewq}[$\star\star$ \iqroles{researcher}]\label{iq:hf:hedging-and-inventory-in-practice:2}
How would you estimate the betas used to hedge a book of forty names with the index future, and how would you test the hedge?
\end{interviewq}

\begin{interviewq}[$\star\star$ \iqroles{researcher}]\label{iq:hf:hedging-and-inventory-in-practice:3}
Why does a proportional cost produce a no-trade band, and a fixed cost a jump back to zero?
\end{interviewq}

\begin{interviewq}[$\star\star$ \iqroles{risk}]\label{iq:hf:hedging-and-inventory-in-practice:4}
Three desks each report a small index exposure within limits. What could go wrong at the firm level?
\end{interviewq}

\begin{interviewq}[$\star\star$ \iqroles{developer}]\label{iq:hf:hedging-and-inventory-in-practice:5}
Design the service that computes the firm's \omterm{def:hf:hedging-and-inventory-in-practice:delta}{delta-equivalent inventory} in real time across stocks, ETFs, futures and options.
\end{interviewq}

\begin{interviewq}[$\star\star\star$ \iqroles{researcher}]\label{iq:hf:hedging-and-inventory-in-practice:6}
Derive the half-width of the dead zone and explain how it changes with the horizon of the hedge.
\end{interviewq}
